%%%%%%%%%%%%%%%%
%%% num 12 %%%
%%%%%%%%%%%%%%%%

\Chapter{12}

\Verse{1}
{
	\hE{וַתְּדַבֵּר מִרְיָם וְאַהֲרֹן בְּמֹשֶׁה עַל אֹדוֹת הָאִשָּׁה הַכֻּשִׁית אֲשֶׁר לָקָח כִּי אִשָּׁה כֻשִׁית לָקָח׃}
}
{
	\eE{
		12 And Miriam and Aaron spoke against Moses about the
		Cushite wife whom he had taken—because he had taken a
		Cushite wife.
	}
}
{
%	\footnote{
%		the Cushite wife. Cush is usually understood to be Ethiopia. Its people
%		are identified as descendants of Noah’s son Ham (Gen 10:6–7). On this
%		understanding, Moses has taken an Ethiopian wife in addition to his
%		first, Midianite wife, Zipporah. This has been confused somewhat by the
%		fact that the prophet Habakkuk refers to a place called Cushan in
%		parallel with Midian (Hab 3:7). Some scholars, therefore, have concluded
%		that the Cushite wife is Zipporah herself (so Ibn Ezra). But this latter
%		view does not explain why the text should suddenly refer to her here as
%		a Cushite. Also, the words “because he had taken a Cushite wife”
%		certainly appear to be here in the verse in order to inform us of an
%		essential new fact, but we already knew that he was married to Zipporah
%		(so Rashbam). This story, therefore, is about a second, probably
%		Ethiopian wife.
%	}
}

\Verse{2}
{
	\hE{וַיֹּאמְרוּ הֲרַק אַךְ בְּמֹשֶׁה דִּבֶּר יְהֹוָה הֲלֹא גַּם בָּנוּ דִבֵּר וַיִּשְׁמַע יְהֹוָה׃}
}
{
	\eE{
		And they said, “Has YHWH only just spoken through Moses?
		Hasn’t He also spoken through us?” And YHWH heard.
	}
}
{
%	\footnote{
%		only just spoken. Aaron and Miriam say, “Has YHWH only just spoken
%		through Moses?” “Only just” sounds redundant. The redundant word “only,”
%		though, is Hebrew raq, and the story concludes with YHWH demanding a
%		measure of punishment of Miriam comparable to the shame she would have
%		experienced if her father “had spit in her face,” which is Hebrew
%		y[image]r[image]q y[image]raq (12:14), an intensive form that doubles
%		the verb—and doubly emphasizes the pun. And the pun confirms that
%		Miriam’s punishment is meant to fit the crime: see the comment on v. 10.
%		One might regard this last pun as a relatively minor item of wordplay,
%		but it is in fact part of an elaborate chain of puns that are stitched
%		through this section of Numbers. The chain encompasses this story and
%		the one that precedes it, the story of Moses’ anguished complaint and
%		the appointment of the seventy elders to share his load. The first word
%		of that story is the extremely unusual and uncertain term ’sapsup
%		(spelled with an unvocalized aleph at the beginning because it is
%		preceded by the definite article), meaning a “gathered mass”
%		(classically translated as a “mixed multitude”) among the Israelites.
%		The root letters of this strange word (’SP) then resurface in a showcase
%		of forms through this story and the next one: [image] Through fourteen
%		variations on the root letters, we are brought from the crisis that is
%		generated by the “gathered mass” to resolution when Miriam is “gathered
%		back” (h[image]’[image]s[image]p). Moreover, this chain of paronomasia
%		may serve to stitch these two stories more tightly together and thus
%		suggest a relationship between the seventy elders’ prophesying in
%		Numbers 11 and Aaron’s and Miriam’s complaint in Numbers 12 (“Has YHWH
%		only just spoken through Moses?”). Footnote 12:2. they said, “Has YHWH
%		only just spoken through Moses?” What does this complaint have to do
%		with Moses’ wife?! Answer: nothing. Like someone who has had a bad day
%		and comes home and complains about dinner, Aaron and Miriam do not voice
%		what is really bothering them. They redirect their annoyance at Moses
%		himself. Aaron and Miriam have been identified as prophets (Exod 7:1;
%		15:20). God has spoken to Aaron, and Aaron was part of the group that
%		had a vision of God on Sinai (Exodus 24). And so they question Moses’
%		position: is he the only one close to God? It is common for people to
%		redirect anger, but it has risks in every case, especially when it
%		occurs within a family. This episode dramatizes that danger: Miriam and
%		Aaron bring God into what was a family dispute. This is a mistake! God
%		appears, and they suffer.
%	}
}

\Verse{3}
{
	\hE{וְהָאִישׁ מֹשֶׁה עָנָו מְאֹד מִכֹּל הָאָדָם אֲשֶׁר עַל פְּנֵי הָאֲדָמָה׃}
}
{
	\eE{
		And the man Moses was very humble, more than every human
		who was on the face of the earth. [image]
	}
}
{
%	\footnote{
%		humble, more than every human who was on the face of the earth. Is this
%		the humblest man on earth? But he is the man who has just said things to
%		God like “Why have you done bad to your servant?” and “If this is how
%		you treat me, kill me!” We might say that his closeness to God enables
%		him to speak this way to the deity while still being humble when dealing
%		with other human beings. But then how are we to account for the way he
%		spoke to the Pharaoh before the last plague in Egypt: “All these
%		servants of yours will come down to me, and they’ll bow to me” (Exod
%		11:8)? Perhaps it is precisely his experiences in Egypt and at Sinai
%		that have made him now the humblest of humans. Having seen the extent of
%		God’s power, he feels humble. Having seen his inability to handle the
%		people and their complaints on his own, he feels humbler still. Thus,
%		even when he speaks to God with his audacious closeness, he says, “I’m
%		not able, I, by myself, to carry this entire people, because it’s too
%		heavy for me” (11:14).
%	}
}

\Verse{4}
{
	\hE{וַיֹּאמֶר יְהֹוָה פִּתְאֹם אֶל מֹשֶׁה וְאֶל אַהֲרֹן וְאֶל מִרְיָם צְאוּ שְׁלשְׁתְּכֶם אֶל אֹהֶל מוֹעֵד וַיֵּצְאוּ שְׁלשְׁתָּם׃}
}
{
	\eE{
		And YHWH said suddenly to Moses and to Aaron and to
		Miriam, “Go out, the three of you, to the Tent of
		Meeting.” And the three of them went out.
	}
}
{
}

\Verse{5}
{
	\hE{וַיֵּרֶד יְהֹוָה בְּעַמּוּד עָנָן וַיַּעֲמֹד פֶּתַח הָאֹהֶל וַיִּקְרָא אַהֲרֹן וּמִרְיָם וַיֵּצְאוּ שְׁנֵיהֶם׃}
}
{
	\eE{
		And YHWH came down in a column of cloud and stood at the
		entrance of the tent. And He called, “Aaron and Miriam.”
		And the two of them went out.
	}
}
{
}

\Verse{6}
{
	\hE{וַיֹּאמֶר שִׁמְעוּ נָא דְבָרָי אִם יִהְיֶה נְבִיאֲכֶם יְהֹוָה בַּמַּרְאָה אֵלָיו אֶתְוַדָּע בַּחֲלוֹם אֲדַבֶּר בּוֹ׃}
}
{
	\eE{
		And He said, “Hear my words: If there will be a prophet
		among you, I, YHWH, shall be known to him in a vision; in
		a dream I shall speak through him.
	}
}
{
%	\footnote{
%		in a vision; in a dream. All prophetic experience in the Tanak is
%		understood to be through visions and dreams—except Moses’. The fifteen
%		books of the Hebrew Bible that are named for prophets either identify
%		the prophets’ experiences as visions or else leave the form of the
%		experiences undescribed (Ezek 12:27; 40:2; Hos 12:11; Hab 2:2; Mic 3:6).
%		Many begin by identifying the book’s contents as the prophet’s vision:
%		“The vision of Isaiah” (Isa 1:1; cf. 2 Chr 32:32); “The vision of
%		Obadiah” (Oba 1); “The book of the vision of Nahum” (Nah 1:1); “The
%		words of Amos … which he envisioned” (Amos 1:1); “The word of YHWH that
%		came to Micah … which he envisioned” (Mic 1:1); “The oracle that
%		Habakkuk the prophet envisioned” (Hab 1:1).
%	}
}

\Verse{7}
{
	\hE{לֹא כֵן עַבְדִּי מֹשֶׁה בְּכל בֵּיתִי נֶאֱמָן הוּא׃}
}
{
	\eE{
		Not so is my servant Moses; in all my house he is
		faithful.
	}
}
{
%	\footnote{
%		Not so is my servant Moses. Some experiences are impressive or inspiring
%		or powerful in themselves; and some experiences change a person’s life
%		ever after. The moment at which Moses sees the form of God from the back
%		on Sinai separates him from all other prophets (Exod 33:21–23). It is
%		understood that no other prophet has a standing that is equal to that of
%		Moses.
%	}
}

\Verse{8}
{
	\hE{פֶּה אֶל פֶּה אֲדַבֶּר בּוֹ וּמַרְאֶה וְלֹא בְחִידֹת וּתְמֻנַת יְהֹוָה יַבִּיט וּמַדּוּעַ לֹא יְרֵאתֶם לְדַבֵּר בְּעַבְדִּי בְמֹשֶׁה׃}
}
{
	\eE{
		Mouth to mouth I shall speak through him and vision and
		not in enigmas, and he will see the form of YHWH. And why
		did you not fear to speak against my servant, against
		Moses?”
	}
}
{
}

\Verse{9}
{
	\hE{וַיִּחַר אַף יְהֹוָה בָּם וַיֵּלַךְ׃}
}
{
	\eE{And YHWH’s anger flared against them, and He went. [image]}
}
{
}

\Verse{10}
{
	\hE{וְהֶעָנָן סָר מֵעַל הָאֹהֶל וְהִנֵּה מִרְיָם מְצֹרַעַת כַּשָּׁלֶג וַיִּפֶן אַהֲרֹן אֶל מִרְיָם וְהִנֵּה מְצֹרָעַת׃}
}
{
	\eE{
		And the cloud turned from over the tent; and, here, Miriam
		was leprous, like snow. And Aaron turned to Miriam, and,
		here, she was leprous.
	}
}
{
%	\footnote{
%		leprous, like snow. Miriam is left with a kind of leprosy in which her
%		flesh is devoid of pigment. If we understand correctly that the Cushite
%		wife is Ethiopian, then this is a provocative case of punishment to suit
%		the crime and a powerful biblical statement regarding racism. It is as
%		if God says to Miriam: “You don’t like a woman with dark pigmentation?
%		Then you don’t have to have any pigmentation!” Or: “You don’t like a
%		woman who is black? Try being completely white!”
%	}
}

\Verse{11}
{
	\hE{וַיֹּאמֶר אַהֲרֹן אֶל מֹשֶׁה בִּי אֲדֹנִי אַל נָא תָשֵׁת עָלֵינוּ חַטָּאת אֲשֶׁר נוֹאַלְנוּ וַאֲשֶׁר חָטָאנוּ׃}
}
{
	\eE{
		And Aaron said to Moses, “In me, my lord. Don’t set a sin
		on us, which we did foolishly and which we sinned.
	}
}
{
%	\footnote{
%		my lord. Aaron addressed Moses as “my lord” one other time: in the
%		golden calf episode. On each of the two occasions on which he has acted
%		inappropriately toward Moses, he has ended up having to acknowledge
%		Moses’ status humbly and politely. Footnote 12:11. Don’t set a sin on
%		us. Aaron acknowledges that both he and Miriam are guilty. He says,
%		“Don’t set a sin on us, which we did foolishly and which we sinned.” Why
%		is there no consequence here for Aaron? Why does Miriam suffer while
%		Aaron does not? As I commented in the case of Aaron’s making the golden
%		calf, the answer may be that Aaron is the high priest. In that position
%		he cannot suffer any direct punishment from God, which would both
%		disqualify him and demean the office of high priest. And so what is left
%		is for Aaron to make amends through personal atonement. What is
%		impressive here is that he frames his atonement entirely in terms of
%		Miriam. He acknowledges his own equal culpability, and he pleads with
%		Moses to intercede for his sister.
%	}
}

\Verse{12}
{
	\hE{אַל נָא תְהִי כַּמֵּת אֲשֶׁר בְּצֵאתוֹ מֵרֶחֶם אִמּוֹ וַיֵּאָכֵל חֲצִי בְשָׂרוֹ׃}
}
{
	\eE{
		Let her not be like the dead who, when he comes out of his
		mother’s womb, half of his flesh is eaten up!”
	}
}
{
}

\Verse{13}
{
	\hE{וַיִּצְעַק מֹשֶׁה אֶל יְהֹוָה לֵאמֹר אֵל נָא רְפָא נָא לָהּ׃}
}
{
	\eE{
		And Moses cried out to YHWH, saying, “Oh, {\God}, heal her!”
		{[image]}
	}
}
{
%	\footnote{
%		Oh, God, heal her! The untranslatable Hebrew n[image]’ particle occurs
%		twice in Moses’ words, conveying that this is a fervent plea, something
%		like: “God, please, heal her, please!”
%	}
}

\Verse{14}
{
	\hE{וַיֹּאמֶר יְהֹוָה אֶל מֹשֶׁה וְאָבִיהָ יָרֹק יָרַק בְּפָנֶיהָ הֲלֹא תִכָּלֵם שִׁבְעַת יָמִים תִּסָּגֵר שִׁבְעַת יָמִים מִחוּץ לַמַּחֲנֶה וְאַחַר תֵּאָסֵף׃}
}
{
	\eE{
		And YHWH said to Moses, “And if her father had spit in her
		face, wouldn’t she be humiliated seven days? Let her be
		closed up seven days outside the camp, and after that let
		her be gathered back.”
	}
}
{
}

\Verse{15}
{
	\hE{וַתִּסָּגֵר מִרְיָם מִחוּץ לַמַּחֲנֶה שִׁבְעַת יָמִים וְהָעָם לֹא נָסַע עַד הֵאָסֵף מִרְיָם׃}
}
{
	\eE{
		And Miriam was closed up outside the camp seven days. And
		the people did not travel until Miriam was gathered back,
		and after that the people traveled from Hazeroth. And they
		camped in the Paran wilderness.
	}
}
{
}

\Verse{16}
{
	\hE{וְאַחַר נָסְעוּ הָעָם מֵחֲצֵרוֹת וַיַּחֲנוּ בְּמִדְבַּר פָּארָן׃}
}
{
	\eE{[image]}
}
{
}
